%% ============================================================
%%  SOLUCIÓN — Ejercicio 2.1: Tabla de regresión al estilo AER
%%  Clase 2 — Después de diapositiva "Tabla de regresión al estilo AER"
%%  Compilar: F5 (pdflatex, dos veces)
%% ============================================================

\documentclass[12pt, a4paper]{article}
\usepackage[utf8]{inputenc}
\usepackage[T1]{fontenc}
\usepackage[spanish]{babel}
\usepackage{amsmath, amssymb}
\usepackage{booktabs}

\title{Ecuación de Mincer: Análisis de regresión}
\author{Tu Nombre}
\date{\today}

\begin{document}
\maketitle

\section{Retornos a la educación}

La Tabla \ref{tab:mincer} presenta los resultados de estimar la ecuación de 
Mincer bajo tres especificaciones distintas. En la columna (1) estimamos solo 
el efecto de educación. En (2) agregamos experiencia potencial y su cuadrado. 
En (3) incluimos un indicador de género.

\begin{table}[htbp]
  \centering
  \caption{Ecuación de Mincer: retornos a la educación}
  \label{tab:mincer}
  \begin{tabular}{lccc}
    \toprule
                         & (1)        & (2)        & (3)       \\
                         & \multicolumn{3}{c}{Variable dep.: $\ln(salario)$} \\
    \cmidrule(lr){2-4}
    \midrule
    Educación            & 0.082***   & 0.071***   & 0.069***  \\
                         & (0.008)    & (0.009)    & (0.010)   \\[5pt]
    Experiencia          &            & 0.034**    & 0.031**   \\
                         &            & (0.012)    & (0.013)   \\[5pt]
    Experiencia$^2$      &            & $-$0.001** & $-$0.001* \\
                         &            & (0.000)    & (0.000)   \\[5pt]
    Mujer (=1)           &            &            & $-$0.183*** \\
                         &            &            & (0.042)   \\[5pt]
    Constante            & 7.241***   & 7.018***   & 7.312***  \\
                         & (0.098)    & (0.115)    & (0.122)   \\
    \midrule
    Controles sector     & No         & No         & Sí        \\
    Observaciones        & 1.200      & 1.200      & 1.200     \\
    $R^2$ ajustado       & 0.241      & 0.318      & 0.374     \\
    \bottomrule
  \end{tabular}
  \begin{minipage}{\linewidth}
    \footnotesize \textit{Notas:} Errores estándar robustos a heterocedasticidad 
    en paréntesis. *** $p<0.01$, ** $p<0.05$, * $p<0.1$. La especificación 
    (3) incluye controles por sector económico (8 categorías). 
    \textit{Fuente:} EPH 2023, elaboración propia.
  \end{minipage}
\end{table}

\section{Respuestas a las preguntas}

\subsection*{¿Qué hace \texttt{\textbackslash cmidrule(lr)\{2-4\}}?}

Este comando dibuja una línea horizontal parcial (no completa) que abarca 
columnas 2 a 4. Los parámetros \texttt{(lr)} agregan pequeñas extensiones 
a izquierda (``l'') y derecha (``r'') para que la línea se vea mejor 
tipográficamente. Sin los parámetros, la línea estaría perfectamente centrada 
en las celdas.

\subsection*{¿Cómo obtener el signo negativo tipográfico correcto?}

La diferencia es:
\begin{itemize}
  \item \textbf{Incorrecto:} \texttt{-0.001} produce ``-0.001'' (guion corto)
  \item \textbf{Correcto:} \texttt{\$-\$0.001} o simplemente \texttt{\$-0.001\$} 
        en modo matemático produce ``$-0.001$'' (signo menos verdadero)
\end{itemize}

En la tabla, usamos la notación \verb|-0.001| que se interpreta como 
``menos'' porque estamos en modo texto pero \LaTeX{} es lo suficientemente 
inteligente como para renderizar el signo correcto. Una forma explícita es 
usar modo matemático: \verb|$-$0.001|.

\section*{Interpretación económica}

\begin{enumerate}
  \item \textbf{Retorno a la educación:} Un año adicional de educación aumenta 
        el salario en aproximadamente 6.9\% (especificación 3).
  \item \textbf{Retorno a la experiencia:} Cada año de experiencia aumenta el 
        salario en 3.1\%, pero el efecto decrece con los años de experiencia 
        (coeficiente negativo en experiencia$^2$).
  \item \textbf{Brecha de género:} Las mujeres ganan 18.3\% menos que los hombres, 
        controlando por educación y experiencia. Esta brecha podría reflejar 
        discriminación, segregación ocupacional, o diferencias en negociación de salarios.
\end{enumerate}

\end{document}
